(1) 弹簧伸长增量 \(y-a\theta\)、\(y+a\theta\)，\(T=(m\dot y^2+I\dot\theta^2)/2\)，\(V=[k_1(y-a\theta)^2+k_2(y+a\theta)^2]/2\)。质量矩阵 \(\mathsf M=\operatorname{diag}(m,I)\)，刚度矩阵 \(\mathsf K=\begin{pmatrix}K&aD\\aD&a^2K\end{pmatrix}\)，\(K=k_1+k_2,D=k_2-k_1\)。
(2) \(\det(\mathsf K-\omega^2\mathsf M)=0\) 给出 \(mI\omega^4-K(I+ma^2)\omega^2+4a^2k_1k_2=0\)，故 \(\omega_\pm^2=\frac12[K/m+a^2K/I\pm\sqrt{(K/m-a^2K/I)^2+4a^2D^2/(mI)}]\)。
(3) 因 \(K>0\)、\(\det\mathsf K=4a^2k_1k_2>0\)，刚度正定，故两频率均实且正。\(k_1=k_2=k\) 时 \(D=0\)，纯平移频率 \(\omega_y^2=2k/m\)，纯转动频率 \(\omega_\theta^2=2ka^2/I\)，模态质量正交。
(4) 非对称耦合强度为 \(aD/\sqrt{mI}\)；原频率平方相遇时仍有劈裂 \(2|aD|/\sqrt{mI}\)，因此避免交叉。对 \(\mathsf M\ddot{\mathbf q}+\mathsf K\mathbf q=0\) 左乘 \(\dot{\mathbf q}^{\mathsf T}\)，严格得到 \(d(T+V)/dt=0\)。
